\scanpage{9}
[\uline{NB}. $\prod_{k'/k}G_{k'}$ est de t. m.] donc grâce à
$H^{n+2}=0$ sur les groupes de t. m. on voit que
\diagram{\begin{tikzcd}
H^{n+1}(\prod_{k'/k}G_{k'})\arrow[r]\arrow[d,phantom,"\unsure{\simeq}" description,sloped]&H^{n+1}(G)\\
H^{n+1}(G_{k'})&
\end{tikzcd}}
est surjectif

donc pour vérifier que $H^{n+1}(G)=0$, il suffit
de le faire pour $H^{n+1}(G_{k'})$. Or
\unsure{nous} \unsure{voici|ramenons} au cas où $G$ est ``\unsure{splitté}''
par une $p$-extension, \uline{ce qui implique}
\uline{que} $G$ \uline{se dévisse en groupes isomorphes}
\uline{à} $\mu_p$. Donc on est ramené à vérifier
que $H^{n+1}(\mu_p)=0$, ce qui se voit sur
la suite exacte
\diagram{\begin{tikzcd}
0\arrow[r]&\mu_p\arrow[r]&G_m\arrow[r]&G_m\arrow[r]&0
\end{tikzcd}}

\uline{Lemme} 3 \deleted{\unsure{Si}} On a
\[
\mathrm{cd}^{\substack{\text{rad}\\\text{add}}}_p(k)=
\begin{cases}0&\text{si }k\text{ parfait}\\1&\text{si }k\text{ imparfait}\end{cases}
\]
En effet, un groupe radiciel additif se dévisse en
groupes $\alpha_p$ \unsure{\ldots}
\mspage{8}

